\documentclass[12pt,letterpaper]{report}
\usepackage{bookman,eulervm}
\usepackage{amsmath,amssymb}
\usepackage{amsthm}
\usepackage[margin=1in]{geometry}
\usepackage{graphicx}

\numberwithin{equation}{chapter}

\newcommand\chap[1]{%
	\chapter*{#1}%
	\addcontentsline{toc}{chapter}{#1}}

\begin{document}
\title{Ordinary Difference Equations}
\author{Eric Conrad}
\date{\today}
\maketitle

\tableofcontents

\chap{Introduction}

An ordinary difference equation is a relation of the following form:
\begin{align}
	f(a_n, a_{n-1}, \dots, a_{n-k}) = 0 \text{.}	\tag{A}
\end{align}
The subscripts are known as indices and typically range over either the natural numbers
or the integers.\footnote{Here we include zero as a natural number.}  If the indices
in the relation are integers that range from $n-k$ through $n$, then the \emph{order}
of the equation is $k$.

The indices are often associated with time.  Hence, conditions on the function $a$ at specific
index values are known as \emph{initial conditions}.

\chapter{The Exponential difference equation}

Recall the following ordinary differential equation from calculus:
\begin{align}
		\frac{dy}{dt} &= cy\text{.} \\
\intertext{The general solution may be found in a number of ways, including separation of
	variable and undetermined coefficients:}
		y &= Ae^t\text{.} \\
\intertext{If we know the value of the function at a given point, we can find a specific
	solution.  For example, if $y(0)=5$, then:}
		y &= 5e^t\text{,} \\
\intertext{or if $y(1)=7$, then:}
		y &= 7e^{t-1}\text{.}
\end{align}

The analogous difference equation is:
\begin{align}
		\frac{y(t+h)-y(t)}{h} &= cy(t)\text{.} \label{eqn:diffForm} \\
\intertext{Normally we have $h=\Delta t=1$, so this reduces to:}
		y(t+1)-y(t) &= cy(t)\text{.} \\
\intertext{But we can combine the expression on the right with the $y(t)$ term on the left:}
		y(t+1) &= (c+1)y(t) \text{.} \\
\intertext{Assuming $c\neq -1$, iterating this expression (or using some other resonable method)
	yields the following solution:}
		y(t) = y(0) (c+1)^t\text{,}
\end{align}
{for all integers $t$.}	

In particular, if $c=1$, we obtain $y=y(0) 2^t$ for integers $t$.  So the difference equation analogue
of $e$ is the number 2.

Equation \ref{eqn:diffForm} is expressed in difference form.  It is often more convenient to write
an ordinary difference equation in shift form -- this is typically done when using index notation.
In shift form, we would typically write:
\begin{align}
	a_t &= ca_{t-1}\text{.} \label{eqn:shiftForm} \\
	\intertext{Assuming $c\neq -1$, iterating this expression yields the following solution:}
	a(t) &= a_0 c^t\text{,}
\end{align}
{for all integers t.}	

It is easy to see that Equations \ref{eqn:diffForm} and \ref{eqn:shiftForm} are closely related.  A solution
to either one can easily be translated into a solution to the other.

\section{Undetermined coefficients}

Now let us solve Equation \ref{eqn:shiftForm} using the method of undetermined coefficients.  Here we need
to know \emph{a priori} that the equation has a particular solution of the form $y=\lambda^n$.  Then:
\begin{align}
	\lambda^n &= c\lambda^{n-1}\text{.} \\
	\intertext{Subtracting the right-hand side and factoring::}
	\lambda^{n-1} (\lambda-c) &= 0\text{.} \\v
	\intertext{To avoid trivialities, we look at $\lambda-c=0$.}
	\lambda &= c\text{.} \\
	\intertext{The difference equation is first order, homogeneous, and linear, so we have one term with
		a coefficient:}
	a_n &= kc^n\text{.} \\
	\intertext{Setting $n=0$:}
	a_0 &= kc^0 = k\text{, so} \\
	a_n &= a_0c^n\text{.}
\end{align}

\chapter{The Fibonacci difference equation}

In a somewhat contrived rabbit breeding example, Leonardo of Pisa (aka Fibonacci) looked at the following
difference equation:
\begin{align}
	f_n &= f_{n-1} + f_{n-2}\text{, subject to} \\
	f_0 &= 0\text{.}
	f_1 &= 1\text{.}
\end{align}

This is a second order linear homogeneous difference equation.  We will use the method of undetermined
coefficients by first finding two \emph{independent} particular solutions.  We try a solution of the
form $y=\lambda^n$:
\begin{align}
	\lambda^n - \lambda^{n-1} - \lambda^{n-2} &= 0 \text{,} \\
	\lambda^{n-2} (\lambda^2 - \lambda - 1) &= 0 \text{.} \\
\intertext{The first factor is nonzero, so our characteristic polynomial is the expression in parentheses:}
	\lambda^2 - \lambda - 1 &= 0 \\
	\lambda &= \frac{1\pm\sqrt{5}}{2} \text{.}
\end{align}

The roots (or \emph{eigenvalues}) are $\phi:=(1+\sqrt{5})/2$ and $\bar{\phi}:=(1-\sqrt{5})/2$.  So the
general solution is:
\begin{align}
	f_n &= A\phi^n + B\bar{\phi}^n\text{.}
\end{align}

Note that $\phi+\bar{\phi}=1$ and $\phi-\bar{\phi}=\sqrt{5}$.  Applying the initial conditions, we know
that:
\begin{align}
	0 = f_0 &= A + B\text{, and} \\
	1 = f_1 &= A\phi + B\bar{\phi}\text{.} \\
	\intertext{From the first equation: $B=-A$.  Substituting into the second:}
	1 &= A\phi - A\bar{\phi} = A\sqrt{5}\text{.} \\
\intertext{Then $A = 1/\sqrt{5}$ and:}
	a_n &= \frac{1}{\sqrt{5}}(\phi^n - \bar{\phi}^n)\text{.}
\end{align}

Some small values are given in Table \ref{tbl:fib}:
\begin{table}[ht]
	\begin{center}
	\begin{tabular}{r|rrrrrrrrrrrrr}
	$n$      & $0$ & $1$ &  $2$ & $3$ & $4$  & $5$ &  $6$ &  $7$ &   $8$ &  $9$ &  $10$
			 & $11$ &  $12$ \\
	\hline
	$f_n$    & $0$ & $1$ &  $1$ & $2$ & $3$  & $5$ &  $8$ & $13$ &  $21$ & $34$ &  $55$
			 & $89$ &  $144$ \\ 
	$f_{-n}$ & $0$ & $1$ & $-1$ & $2$ & $-3$ & $5$ & $-8$ & $13$ & $-21$ & $34$ & $-55$
			 & $89$ & $-144$ \\ 
	\end{tabular}
	\end{center}
	\caption{Fibonacci numbers for small indices}
	\label{tbl:fib}
\end{table}

As we see in the table, the sequence is bilateral.  From $f_{-(n+2)}+f_{-(n+1)}=f_{-n}$, it follows that:
\[
	f_{-(n+2)} = f_{-n} - f_{-(n+1)}\text{.}
\]
Using mathematical induction, it isn't difficult to prove that:
\begin{align}
	f_{-n} = (-1)^{n+1} f_n \text{.}
\end{align}


\end{document}

